\section{Introducción y Justificación}
\label{sec:introduccion}
% En este apartado se debe exponer el problema clínico y metalúrgico. Se analizan los métodos convencionales de conformado plástico (forja, laminación, trefilado) y sus limitaciones mecánicas en implantes personalizados. Es crítico introducir aquí el concepto de módulo elástico (módulo de Young) y el fenómeno del \textit{stress shielding} (protección de carga) que sufren los tejidos óseos cuando se implantan estructuras excesivamente rígidas.\cite{park,ratner}


Los primeros intentos de utilizar titanio en la fabricación de implantes se remontan a finales de la década de 1930, cuando se demostró su tolerancia en fémures de modelos animales \cite{park}. Desde entonces, el titanio y en particular la aleación Ti-6Al-4V se han consolidado como materiales de referencia en sectores de alta exigencia como la ingeniería aeroespacial y biomédica \cite{park,ratner,nguyen2022critical}. Su adopción masiva se justifica por su excelente resistencia a la corrosión, su elevada resistencia específica ---la mayor entre los elementos metálicos--- y su baja densidad ($4{,}5\;\text{g/cm}^3$), notablemente inferior a la del acero inoxidable 316 ($7{,}9\;\text{g/cm}^3$), la aleación CoCrMo ($8{,}3\;\text{g/cm}^3$) o la aleación CoNiCrMo ($9{,}2\;\text{g/cm}^3$) \cite{park}. Sin embargo, el mecanizado y conformado convencional del Ti-6Al-4V presenta desafíos significativos \cite{nguyen2022critical}. Bajo condiciones extremas de trabajo, esta aleación exhibe un mecanismo de deformación inusual caracterizado por el ablandamiento térmico. Esta propiedad, combinada con su baja conductividad térmica y bajo calor específico, provoca temperaturas de corte extremadamente altas que aceleran el desgaste de las herramientas de mecanizado convencionales \cite{nguyen2022critical}.

No obstante, más allá de las dificultades de fabricación, el Ti-6Al-4V macizo presenta una limitación clínica fundamental: su módulo elástico ($100$--$110\;\text{GPa}$) es muy superior al del hueso cortical ($12$--$20\;\text{GPa}$) y al del hueso esponjoso ($<4\;\text{GPa}$) \cite{park,ratner,misiaszek20243d}. Esta disparidad de rigidez provoca el fenómeno conocido como \textit{stress shielding} (apantallamiento de tensiones): el implante soporta la mayor parte de la carga mecánica, lo que reduce la estimulación del tejido óseo adyacente y conduce a su reabsorción progresiva según la ley de Wolff \cite{park,ratner}. Para mitigar este efecto, la estrategia más prometedora consiste en diseñar estructuras porosas o de tipo \textit{lattice} (estructura reticular) que reduzcan el módulo efectivo del implante y, simultáneamente, favorezcan la osteointegración al permitir el crecimiento óseo en su interior \cite{misiaszek20243d,li2024biomimetic,kilina2024influence}.

Como alternativa a los métodos de conformado convencional, la Fabricación Aditiva (AM, por sus siglas en inglés) ha emergido como un proceso transformador que permite materializar estas geometrías porosas complejas \cite{nguyen2022critical}. La AM construye componentes capa por capa a partir de modelos de diseño asistido por ordenador (CAD) utilizando fuentes de calor como láser, haz de electrones o ultrasonidos para fundir y consolidar el material. Este enfoque produce geometrías casi definitivas (\textit{near-net-shape}), reduciendo drásticamente el desperdicio de material y acortando los tiempos de fabricación al omitir etapas tradicionales como el moldeo o el mecanizado intensivo \cite{nguyen2022critical}.

\subsection{Definición de Tecnologías de Fabricación Aditiva}
Para el procesamiento de aleaciones de titanio, la industria se apoya principalmente en dos categorías tecnológicas: la Fusión de Lecho de Polvo (PBF) y la Deposición de Energía Dirigida (DED). La tecnología PBF utiliza una fuente de calor (láser o haz de electrones) para fundir selectivamente las partículas de un lecho de polvo metálico pre-esparcido. En contraste, los sistemas DED funden el material (en formato de polvo o hilo) simultáneamente mientras es depositado a través de una boquilla \cite{neikter2018microstructural,zhai2016microstructure}.

Las fuentes de energía definen las atmósferas de trabajo: las tecnologías basadas en haz de electrones (EBM) operan obligatoriamente en condiciones de alto vacío, mientras que los sistemas láser (SLM, DED) emplean atmósferas de gas inerte. Estas diferencias tecnológicas determinan críticamente la historia térmica del material y, en consecuencia, sus propiedades mecánicas finales~\cite{neikter2018microstructural,zhai2016microstructure}.

\subsection{Conceptos Microestructurales y Termodinámicos}
A diferencia de la fabricación convencional que produce fases equiaxiales $\alpha+\beta$, las piezas construidas mediante AM experimentan una historia térmica muy compleja caracterizada por rápidos gradientes de temperatura y ciclos de recalentamiento \cite{nguyen2022critical}. Durante la deposición, el material supera la temperatura \textit{liquidus} y la temperatura de transición \textit{transus} $\beta$, lo que induce la formación de granos columnares previos de fase $\beta$ que crecen epitaxialmente en dirección perpendicular al baño fundido \cite{zhai2016microstructure}.

El factor determinante para la microestructura final es la velocidad de enfriamiento. En los procesos SLM y DED, las velocidades de enfriamiento superan los $410^\circ\text{C/s}$, lo que provoca una transformación difusional completa hacia una fase martensítica acicular metaestable ($\alpha'$) \cite{nguyen2022critical}. Por el contrario, la tecnología EBM mantiene una temperatura de cámara muy elevada (entre $600^\circ\text{C}$ y $760^\circ\text{C}$), lo que actúa como un tratamiento térmico \textit{in situ} que descompone la fase $\alpha'$ para formar una microestructura bifásica $\alpha+\beta$ con entramados de Widmanstätten. Esta distinción explica por qué las piezas fabricadas por SLM y DED exhiben una mayor resistencia a la fluencia pero una menor ductilidad en comparación con las procesadas por EBM \cite{neikter2018microstructural,zhai2016microstructure}.

\subsection{Defectología y Tratamientos de Post-procesado}
A pesar de sus ventajas geométricas, los procesos AM son susceptibles a la formación de defectos internos que limitan la vida a fatiga del implante \cite{zhai2016microstructure}. Se distinguen principalmente dos tipos de porosidades: los poros de gas, de geometría esférica, causados por gases atrapados durante la solidificación; y los defectos por falta de fusión, de geometría irregular y bordes afilados, originados por una fusión incompleta del polvo metálico entre capas \cite{nguyen2022critical}.

Estos poros actúan como concentradores de tensiones e iniciadores de grietas. Para mitigar estos defectos y homogeneizar el comportamiento mecánico, es práctica habitual someter las piezas de Ti-6Al-4V a tratamientos térmicos de post-procesado \cite{zhai2016microstructure}. El Prensado Isostático en Caliente (HIP) destaca como el método más efectivo para cerrar las porosidades internas y densificar la microestructura, mejorando significativamente tanto la ductilidad como el límite de fatiga de los componentes fabricados \cite{nguyen2022critical}.